\title{The Era of 1-bit LLMs: All Large Language Models are in 1.58 Bits}

\begin{abstract}
Emerging studies like BitNet~\cite{bitnet} are ushering in the era of 1-bit Large Language Models (LLMs). In this paper, we present a new 1-bit LLM architecture, termed \textbf{\text{BitNet b1.58}}, wherein every parameter of the model adopts ternary values \{-1, 0, 1\}. This configuration achieves equivalency with full-precision (e.g., FP16 or BF16) Transformer LLMs of comparable model capacity and trained on the same amount of tokens in both perplexity and downstream task outcomes, while offering substantial advantages in latency, memory footprint, throughput, and power efficiency. Notably, the 1.58-bit LLM introduces a novel scaling law and a methodology for training future generations of LLMs that optimize for both high performance and resource efficiency. This direction also supports a new computational paradigm and paves the way for hardware design tailored specifically for 1-bit LLM architectures.
\end{abstract}

\section{The Era of 1-bit LLMs}

The AI community has recently witnessed an explosive increase in the scale and performance of Large Language Models (LLMs). Despite achieving outstanding results across various natural language processing applications, the growing size of these models has introduced obstacles related to their practical deployment, as well as concerns regarding economic and environmental sustainability due to their substantial energy requirements.

To mitigate such issues, post-training quantization has been proposed as a strategy to convert LLMs into low-bit representations for inference~\cite{smoothquant, gptq, quip, quip_sharp}. By reducing the precision of model parameters and activations, this approach delivers significant savings in memory usage and computational overhead. The movement towards lower precision has evolved from 16-bit formats down to 4-bit alternatives~\cite{gptq, awq}. Nevertheless, even though post-training quantization is prevalent in industry applications, it is not an optimal solution.

Recent advancements in 1-bit neural architectures, typified by BitNet~\cite{bitnet}, signal an effective path toward minimizing the computational expenses associated with LLMs without sacrificing efficacy. Traditional LLMs operate primarily with 16-bit floating point formats (such as FP16 or BF16), and their computational intensity is dominated by matrix multiplication. Most of the computational burden thus stems from floating-point additions and multiplications. By comparison, BitNet performs matrix multiplications exclusively using integer additions, which results in significantly lower energy consumption for LLM workloads. Since power often constrains computational throughput in modern hardware, this reduction in energy demand directly supports faster operations.

Beyond raw computation, model inference requires frequent movement of parameters from DRAM into the local accelerator memory (e.g., SRAM), an operation that can be particularly resource-intensive. Some solutions propose scaling up SRAM to boost data throughput, though this comes at a steep financial cost relative to DRAM. One-bit LLMs are inherently more memory-efficient than their full-precision counterparts, decreasing requirements in both storage size and memory bandwidth. This reduced memory footprint cuts both the financial and temporal costs of loading model weights from DRAM, thus accelerating and streamlining inference processes.

In this study, we present a noteworthy 1-bit LLM extension named \textbf{\text{BitNet b1.58}}, in which each model parameter can assume one of three values: \{-1, 0, 1\}. By incorporating 0 alongside the traditional binary states, the bit representation per parameter effectively increases to 1.58 bits. \text{BitNet b1.58} continues to deliver the original strengths of BitNet, including its distinct computational approach that largely eliminates the need for multiplication during matrix operations, thereby enabling further optimization. The model maintains the energy efficiency of its predecessor and exhibits superior memory usage, computational throughput, and reduced latency in comparison with standard FP16 LLM benchmarks. Beyond these benefits, \text{BitNet b1.58} introduces two key improvements: first, it enhances representational power by allowing for explicit feature selection through zero-valued weights, a property that can considerably boost the accuracy of 1-bit LLMs; second, empirical results demonstrate that, starting at the 3B parameter scale, \text{BitNet b1.58} matches the performance of full-precision (FP16) models in terms of both perplexity and downstream task metrics, provided identical training configurations (such as model size and total training tokens).

\section{BitNet b1.58}

\text{BitNet b1.58} builds on the BitNet framework, characterized as a Transformer model where the conventional \emph{nn.Linear} layers are substituted by \emph{BitLinear} layers. This variant is initialized from scratch and operates with weights quantized to 1.58 bits and activations at 8 bits. Compared to the baseline BitNet, this version incorporates several revisions, outlined below.

\paragraph{Quantization Function.} To enforce weight values strictly within \{-1, 0, +1\}, we utilize an \emph{absmean} quantization mechanism. In this process, the weight matrix is normalized by its mean absolute value, following which each element is discretized to the nearest member of \{-1, 0, +1\} via rounding:

\begin{equation}
    \widetilde{W} = \mathrm{RoundClip}\left(\frac{W}{\gamma + \epsilon}, -1, 1\right),
\end{equation}

\begin{equation}
    \mathrm{RoundClip}(x, a, b) = \max(a, \min(b, \mathrm{round}(x))),
\end{equation}

\begin{equation}
    \gamma = \frac{1}{nm}\sum_{ij} |W_{ij}|.
\end{equation}

For activations, the quantization methodology mirrors that of the original BitNet. However, we omit scaling activations to the interval $[0, Q_b]$ before passing them through non-linearities. Instead, each token's activation is scaled to the interval $[-Q_b, Q_b]$, thereby eliminating the need for zero-point quantization. This alteration not only streamlines implementation and facilitates system-level optimizations, but in our empirical analysis, it incurs minimal impact on model performance.

\paragraph{LLaMA-alike Components.}

LLaMA~\cite{llama, llama2} has established itself as the standard backbone architecture within the open-source LLM community. To maximize compatibility and adoption, our implementation of \text{BitNet b1.58} incorporates several LLaMA-inspired architectural features. Notably, it employs RMSNorm~\cite{rmsnorm}, SwiGLU~\cite{swiglu}, rotary embeddings~\cite{rope}, and omits all bias terms. This alignment ensures that \text{BitNet b1.58} can be seamlessly used with widely adopted open-source frameworks such as Huggingface, vLLM~\cite{vllm}, and llama.cpp\footnote{\url{https://github.com/ggerganov/llama.cpp}}, requiring minimal integration overhead.

\section{Results}

We conducted a comparative analysis of \text{BitNet b1.58} and a reproduced FP16 \text{LLaMA LLM} across multiple model scales. For consistency and fairness in benchmarking, all models underwent pre-training using 100 billion tokens from the RedPajama dataset~\cite{redpajama}. Their zero-shot capabilities were assessed on several language benchmarks, specifically ARC-Easy~\cite{arc}, ARC-Challenge~\cite{arc}, Hellaswag~\cite{hellaswag}, Winogrande~\cite{winoGrande}, PIQA~\cite{piqa}, OpenbookQA~\cite{openbookqa}, and BoolQ~\cite{boolq}. Additionally, we presented validation perplexity outcomes on the WikiText2~\cite{wikitext2} and C4~\cite{c4} corpora.

To quantify computational efficiency, we measured both runtime GPU memory utilization and inference latency for \text{LLaMA LLM} and \text{BitNet b1.58}. Evaluations were executed using the FasterTransformer framework\footnote{\url{https://github.com/NVIDIA/FasterTransformer}}, a toolset specifically engineered for optimizing LLM inference speed on GPUs. The 2-bit inference kernel implemented in Ladder~\cite{ladder} was integrated into \text{BitNet b1.58}. We report the elapsed time per output token, as this dominates inference computational cost.

Results summarized in Table~\ref{tab:ppl} reflect both perplexity metrics and resource requirements for \text{BitNet b1.58} and \text{LLaMA LLM}. The findings indicate that, from the 3B parameter scale, \text{BitNet b1.58} achieves perplexity comparable to the full-precision \text{LLaMA LLM}, while delivering 2.71$\times$ speedup and reducing GPU memory usage by a factor of 3.55. Notably, the 3.9B variant of \text{BitNet b1.58} is 2.4$\times$ faster, requires 3.32$\times$ less memory, and outperforms the 3B \text{LLaMA LLM} in accuracy.

Detailed zero-shot evaluation results are shown in Table~\ref{tab:zero-shot}. The assessments were performed following the \emph{lm-evaluation-harness} protocol\footnote{\url{https://github.com/EleutherAI/lm-evaluation-harness}}. These results demonstrate that the difference in zero-shot accuracy between \text{BitNet b1.58} and \text{LLaMA LLM} diminishes as model size increases. Importantly, starting from the 3B configuration, \text{BitNet b1.58} matches the full-precision model on these tasks. Consistent with perplexity trends, the 3.9B \text{BitNet b1.58} instance achieves superior downstream task performance compared to the 3B \text{LLaMA LLM}, with notable gains in memory efficiency and inference speed. This establishes \text{BitNet b1.58} as a Pareto-optimal enhancement relative to leading LLM models.

\paragraph{Memory and Latency}

We extended our experiments to larger model sizes, evaluating versions with 7B, 13B, and 70B parameters, and analyzed their associated computational costs. As depicted in Figure~\ref{fig:latency-memory-energy}, both latency and memory requirements were examined, revealing that the performance acceleration becomes more pronounced with increasing model scale. Notably, the 70B variant of \text{BitNet b1.58} achieved a 4.1-fold speed advantage over the \text{LLaMA LLM} baseline. This improvement arises because the computational demand of \emph{nn.Linear} escalates with model size. Similarly, memory usage demonstrates a comparable pattern; although the embedding layer remains in full precision, its share of total memory diminishes for larger networks. It is important to note that these latency and memory results utilized a 2-bit kernel, suggesting that additional cost reductions may be feasible with further optimizations.

\paragraph{Energy}

We further estimated the energy usage associated with arithmetic operations for both \text{BitNet b1.58} and \text{LLaMA LLM}. Since matrix multiplication dominates the computational burden in LLMs, our analysis centers on this aspect. As shown in Figure~\ref{fig:energy}, the majority of operations in \text{BitNet b1.58} involve INT8 additions, whereas \text{LLaMA LLM} relies mainly on FP16 additions and multiplications. Citing the energy modeling methodology from~\cite{energycost, pokebnn}, the results indicate that \text{BitNet b1.58} requires 71.4 times less energy for matrix multiplication arithmetic on 7nm hardware. We also measured the comprehensive energy costs for full inference runs with 512 input tokens, observing that \text{BitNet b1.58} offers progressively greater energy efficiency as the model size increases, compared to the FP16 \text{LLaMA LLM} reference. This trend occurs because the proportion of computation spent on \emph{nn.Linear} grows with model scale, while the relative energy used by other layers diminishes for larger architectures.

\paragraph{Throughput}

We assessed the throughput of \text{BitNet b1.58} against that of the 70B-parameter \text{LLaMA LLM}, employing two 80GB A100 GPUs and leveraging pipeline parallelism~\cite{gpipe} to facilitate execution of the larger \text{LLaMA LLM} model. The batch size was incrementally raised until we encountered the GPU memory ceiling, while the sequence length remained fixed at 512. As displayed in Table~\ref{tab:throughput}, the 70B \text{BitNet b1.58} accommodates batch sizes as much as 11 times larger than those possible with \text{LLaMA LLM}, resulting in an overall throughput that is 8.9 times higher.

\textbf{\text{BitNet b1.58} introduces a novel scaling paradigm in terms of the interplay between model capability and inference cost.} For illustration, Figures~\ref{fig:latency-memory-energy} and~\ref{fig:energy} reveal equivalencies between various model sizes using 1.58-bit precision and their FP16-based (16-bit) counterparts:

\begin{itemize}
    \item The 13B parameter BitNet b1.58 surpasses the 3B FP16 LLM concerning latency, memory footprint, and energy demand.
    \item The 30B BitNet b1.58 is more performant and efficient—across latency, memory, and energy—than the 7B FP16 LLM.
    \item A 70B BitNet b1.58 demonstrates lower latency, memory use, and energy consumption relative to a 13B FP16 LLM.
\end{itemize}

\paragraph{Training with 2T Tokens}

The scale of training data, specifically the token count, is a key determinant for the performance of LLMs. To examine the token scalability of \text{BitNet b1.58}, we trained it using 2T tokens according to the data pipeline prescribed by StableLM-3B~\cite{StableLM-3B-4E1T}, a leading open-source 3B parameter model. Performance was evaluated on a composite benchmark featuring Winogrande~\cite{winoGrande}, PIQA~\cite{piqa}, SciQ~\cite{sciq}, LAMBADA~\cite{lambada}, and ARC-easy~\cite{arc}. Zero-shot accuracy results, summarized in Table~\ref{tab:2t}, were averaged for metrics combining accuracy and normalized accuracy. StableLM 3b results at the 2T token scale were directly sourced from its official report. The experimental outcomes reveal that \text{BitNet b1.58} outperforms on all benchmarked tasks, suggesting that LLMs with 1.58-bit precision can deliver robust generalization.

\section{Discussion and Future Work}

\textbf{1-bit Mixture-of-Experts (MoE) LLMs}

The Mixture-of-Experts (MoE) architecture has established itself as a highly efficient solution for large language models (LLMs) in terms of computational costs. However, substantial memory usage and intensive inter-chip communication still pose significant barriers to practical adoption. These limitations may be mitigated by employing 1.58-bit LLMs. The smaller memory demands associated with 1.58-bit LLMs lower the hardware requirements for MoE deployment. Additionally, the burden of transferring activations across chips is alleviated. Ideally, eliminating this overhead is possible if the entire model fits within the memory of a single chip.

\textbf{Native Support of Long Sequence in LLMs}

Processing lengthy sequences has become a pivotal requirement for modern LLM applications. A key obstacle in supporting long sequence inference is the substantial memory needed for KV caches. \text{BitNet b1.58} marks a substantial advance toward more effective long-sequence handling, reducing activation storage from 16 bits to 8 bits and, thus, enabling a doubling of context length without additional resource consumption. Further, activations may be compressed to 4 bits or less without loss, particularly in 1.58-bit LLMs; this avenue is reserved for future investigation.

\textbf{LLMs on Edge and Mobile}

Deploying 1.58-bit LLMs can significantly enhance language model operation on edge and mobile hardware, where constraints on memory and processing power frequently hinder LLM deployment. By lowering both memory and power requirements, 1.58-bit LLMs make possible a breadth of new use cases for LLMs on resource-limited platforms. This innovation stands to expand the utility of edge and mobile devices, introducing applications that were not feasible before. Furthermore, 1.58-bit LLMs offer improved compatibility with CPUs, the prevalent processors in such environments, thus enabling efficient execution of \text{BitNet b1.58} and boosting overall device performance and capacity.

\textbf{New Hardware for 1-bit LLMs}

Emerging technologies such as~Groq\footnote{\url{https://groq.com/}} have illustrated both the viability and the promise of creating custom hardware—like LPUs—tailored for LLMs. Building on this, there is an opportunity and a need for the development of novel hardware and system designs specifically optimized for 1-bit LLMs, taking full advantage of the computational models introduced by BitNet~\cite{bitnet}.

\end{document}